\documentclass[10pt,a4paper]{amsart}
\usepackage[margin=19mm]{geometry}
\usepackage{amsmath,amssymb,mathtools}
\usepackage[scr=rsfs,cal=euler]{mathalfa}
\usepackage{xfrac}
\usepackage[unicode,hidelinks,bookmarks=false]{hyperref}
\newcommand{\ie}{i.e.\ }
\newcommand{\cf}{cf.\ }
\newcommand{\resp}{respectively\ }
\newcommand{\Subsection}{\subsection}
\providecommand{\nobreakdash}{\nobreak\hbox{-}}
\setlength{\emergencystretch}{2em}
\multlinegap0pt
\usepackage{bm}
\usepackage{bbm}
\usepackage{dsfont}
\usepackage{xspace}
\usepackage{cancel}
\usepackage{mathtools}
\usepackage{amsthm}
\makeatletter
\@ifundefined{theorem}{\newtheorem{theorem}{Theorem}}{}
\@ifundefined{lemma}{\newtheorem{lemma}[theorem]{Lemma}}{}
\makeatother
\title[Average Precision at Cutoff k under Random Rankings: Expecta]{Average Precision at Cutoff k under Random Rankings: Expectation and Variance}
\author{Tetiana Manzhos, Tetiana Ianevych, Olga Melnyk}
\date{Source: arXiv:2511.02571v1 (CC BY 4.0)}
\begin{document}
\maketitle

\noindent\emph{Source and licence.} Average Precision at Cutoff k under Random Rankings: Expectation and Variance. Version arXiv:2511.02571v1 (CC BY 4.0), distributed under the Creative Commons Attribution 4.0 International licence, as recorded in the arXiv licence metadata for this exact version and shown on the abstract page https://arxiv.org/abs/2511.02571v1. Full rights notice: licenses/arxiv-2511.02571-CC-BY-4.0.txt.\par\smallskip

\noindent\textit{Editorial scope.} Counted unit: the Theorem whose source label is \texttt{th2wor} in the article (source0.tex lines 219--239, source label \texttt{th2wor}), delivered together with the article's own complete original proof of it (source0.tex lines 266--397, one \texttt{proof} environment of 132 lines). No mathematics is written, repaired or re-derived: statement and proof are the article's own text line for line. The proof is detached from the statement in the source: the article states the result at lines 219--239 and prints its proof at lines 266--397, and the proof environment's own header names the statement, so the association is the article's own. Required context retained uncounted: lem1 (lines 242--251). Every definition, display and label of those lines is retained unchanged. Omitted from this package and not redistributed: 1--218, 240--241, 252--265, 398--686. Nothing inside a retained excerpt is omitted or altered: every retained body is delivered exactly as the article's lines are written, so no omission receipt of any kind accompanies this package. Citations: the retained excerpts carry no citation command at all, so no bibliography entry of the article is transcribed and no reference is redistributed. Reference adaptation: none was needed and none was made; every reference inside the delivered unit resolves inside it, so no unresolved reference or citation is produced. Printed slips kept literal: no printed word, symbol, hypothesis or number of the counted statement or of its proof was altered. Layout and class: the article's own class is \texttt{tpms-l} with packages including graphicx; the wrapper is the lane's pinned \texttt{amsart} editorial wrapper at 10pt on A4, which restates the theorem-like environments the retained text needs and adds the article's own macro definitions that the retained text needs (none), with extra wrapper packages (bm, bbm, dsfont, xspace, cancel, mathtools). The delivered numbering is the unit's own. Assets: no retained span contains a figure environment, a graphics-inclusion command, a PDF-page inclusion, a table or a TikZ picture; no image file is copied into this package, the package declares no asset dependency and no image file is redistributed with it. Licence: version arXiv:2511.02571v1 of this article is distributed under the Creative Commons Attribution 4.0 International licence, as recorded in the arXiv licence metadata for this exact version and shown on the abstract page https://arxiv.org/abs/2511.02571v1. A full rights notice is delivered at licenses/arxiv-2511.02571-CC-BY-4.0.txt. Characters: every retained excerpt is pure ASCII, so the delivered engine, XeLaTeX with its default Latin Modern text font, reports no missing character.
\begin{lemma}\label{lem1}
For any $k\geq 1$
\begin{itemize}
\item[(i)]  $\sum_{i=1}^{k-1}\sum_{l=i+1}^{k}\frac{1}{i}  =  k(H_k-1);$
\item[(ii)] $\sum_{i=1}^{k-1} \sum_{l=i+1}^{k} \frac{1}{l}  = k-\sum_{l=1}^{k}\frac{1}{l}=k-H_k;$
\item[(iii)] $\sum_{i=1}^{k-1} \sum_{l=i+1}^{k} \frac{1}{i\cdot l}  =  \frac{1}{2}\bigg(H_k^2-H_k^{\left(2\right)}\bigg),$   
\end{itemize}
where  $H_k=\sum_{i=1}^{k} \frac{1}{i}$ is a standard harmonic number; 
      $H_k^{\left(2\right)}=\sum_{i=1}^{k} \frac{1}{i^2}$ is the $k$-th partial sum of the 2nd order harmonic series.
\end{lemma}

\begin{theorem} \label{th2wor} If we consider the offline evaluation procedure, that is, the random appearance of a fixed number $m$ of relevant items corresponds to sampling $m$ items from $N$ with equal probabilities and without replacement, then 
\begin{equation*}
\begin{split} \operatorname{Var}_{WOR}\big(AP@k\big)=\frac{1}{M^2} \frac{m}{N}\Bigg[& k(C+2(E-F) + (k-1)G) + H_k(B-2(E-kF))+\\ 
& + H_k^2 \cdot D + H_k^{\left(2\right)}(A-D) \Bigg],
\end{split}
\end{equation*}
where 
\begin{align*}
 A & =  1- \frac{m}{N}-\frac{m-1}{N-1} \left(3-2\frac{m-2}{N-2}-\frac{m}{N}\left(2-\frac{m-1}{N-1}\right)\right); \\
 B & =  \frac{m-1}{N-1}\left(3\left(1-\frac{m-2}{N-2}\right)-2\frac{m}{N}\left(1-\frac{m-1}{N-1}\right)\right);\\
C & = \frac{m-1}{N-1}\left(\frac{m-2}{N-2}-\frac{m(m-1)}{N(N-1)}\right);\\
D & = \frac{m-1}{N-1}\left(2-5\frac{m-2}{N-2}+3\frac{(m-2)(m-3)}{(N-2)(N-3)}\right)-\frac{m}{N}\left(1-\frac{m-1}{N-1}\right)^2;\\
E& = \frac{m-1}{N-1}\left(3\frac{m-2}{N-2}\left(1-\frac{m-3}{N-3}\right)-\frac{m}{N}\left(1-\frac{m-1}{N-1}\right)\right);\\
F &= \frac{m-1}{N-1}\left(\frac{m-2}{N-2}\left(1-\frac{m-3}{N-3}\right)-\frac{m}{N}\left(1-\frac{m-1}{N-1}\right)\right); 
\end{align*}
 \begin{align*}
G & = \frac{m-1}{N-1}\left(\frac{(m-2)(m-3)}{(N-2)(N-3)}-\frac{m}{N}\frac{m-1}{N-1}\right);\\
H_{k} & = \sum_{i=1}^k \frac{1}{i}; \quad H^{(2)}_{k} = \sum_{i=1}^k \frac{1}{i^2};\\
M & = \min(m,k). 
\end{align*}
\end{theorem}

\begin{proof}[Proof of Theorem~\ref {th2wor}]
Since $\operatorname{Var}(AP@k) = \operatorname{Var}\left(\frac{1}{k} \sum_{i=1}^k P@i \cdot I_i\right) = \frac{1}{k^2} \sum_{i=1}^k \sum_{l=1}^k c_{il}$, we need to calculate  all  the coefficient $c_{il}$ taking into account that
 for every $i\neq l$
$c_{il} = c_{li} = \operatorname{cov}(P@i \cdot I_i, P@l \cdot I_l)$
and  for all $i$
$c_{ii} = \operatorname{Var}(P@i \cdot I_i)$.

Let us calculate them step by step, starting from the coefficients $c_{ii}$, $i=1,2,...,k.$ Since in this case the random variables $I_1,..., I_n$ are not independent, it is worth applying formulas that use conditional expectations. In particular, for two random variables $X$ and $Y$, we can calculate the variance by the formula:
$$\operatorname{Var}(X)=\operatorname{E}\big(\operatorname{Var}(X\mid Y)\big)+Var\big(\operatorname{E}(X\mid Y)\big).$$
So,
\begin{equation}\label{t1t2}
c_{ii}=\operatorname{Var}\big(P@i\cdot I_i\big)=\operatorname{E}\big(\operatorname{Var}(P@i\cdot I_i\mid I_i\big)+\operatorname{Var} \big(\operatorname{E}(P@i\cdot I_i\mid I_i) \big)=t_1+t_2.
\end{equation}
If $I_i=0,$ then $P@i\cdot I_i=0,$ and $\operatorname{E}\big(P@i\cdot I_i\mid I_i=0\big)=0, \operatorname{Var}\big(P@i\cdot I_i\mid I_i=0\big)=0.$ \\
If $I_i=1,$ then $\operatorname{E}\big(P@i\cdot I_i\mid I_i=1\big)=\frac{1}{i} \left((i-1)\frac{m-1}{N-1}+1 \right)$ and 
$$\operatorname{Var}\big(P@i\cdot I_i\mid I_i=1\big)=\operatorname{E}\left(\big(P@i\cdot I_i\big)^2\mid I_i=1\right)-\big(\operatorname{E}(P@i\cdot I_i\mid I_i=1)\big)^2.$$

First of all
\begin{equation*}
\begin{split}
\operatorname{E}\big((P@i\cdot I_i)^2\mid I_i=1\big)&=\operatorname{E}\left(\left(\frac{1}{i}\left(\sum_{j=1}^{i-1} I_jI_i+I_i^2\right)\right)^2\mid I_i=1\right)=\\
%=\frac{1}{i^2}\operatorname{E}\left(\left(1+\sum_{j=1}^{i-1} %I_j\right)^2\mid I_i=1\right)=\\
=&\frac{1}{i^2}\operatorname{E}\left(1+2\sum_{j=1}^{i-1} I_j
+\left(\sum_{j=1}^{i-1} I_j\right)^2\mid I_i=1\right)=\\
%\end{gather*}
%\begin{gather*}
=&\frac{1}{i^2}\operatorname{E}\left(1+2\sum_{j=1}^{i-1} I_j
+\sum_{j=1}^{i-1} I_j^2+\sum \sum_{j\neq r}I_j I_r\mid I_i=1\right)=\\
=&\frac{1}{i^2}\left(1+3\sum_{j=1}^{i-1} \operatorname{E}\big(I_j \mid I_i=1\big)+
\sum \sum_{j\neq r}\operatorname{E}\big(I_j I_r\mid I_i=1\big)\right)=\\
=&\frac{1}{i^2}\left(1+3(i-1)\cdot \frac{m-1}{N-1}+
(i-1)(i-2)\frac{(m-1)(m-2)}{(N-1)(N-2)}\right),
\end{split}
\end{equation*}
as $\operatorname{E}\big(I_j\mid I_i=1\big)=\frac{C_{N-2}^{m-2}}{C_{N-1}^{m-1}}=\frac{m-1}{N-1}$ and $\operatorname{E}(I_j I_r \mid I_i=1)=\frac{C_{N-3}^{m-3}}{C_{N-1}^{m-1}}=\frac{(m-1)(m-2)}{(N-1)(N-2)}$.

Then \begin{equation*}
\begin{split} 
\operatorname{Var}(P@i\cdot I_i \mid I_i=1)&=\frac{1}{i^2}\left(1+3(i-1)\cdot \frac{m-1}{N-1}+
(i-1)(i-2)\frac{(m-1)(m-2)}{(N-1)(N-2)}\right)-\\
&-\frac{1}{i^2}\left((i-1)\frac{m-1}{N-1}+1 \right)^2
\end{split} 
\end{equation*}
and thus we obtain the first term in \eqref{t1t2}: 
\begin{equation*}
\begin{split} &t_1=\operatorname{E}\big(\operatorname{Var}(P@i\cdot I_i \mid I_i)\big)=\operatorname{Var}\big(P@i\cdot I_i \mid I_i=1\big)\cdot \frac{m}{N}=\\
&=\frac{m}{N} \cdot \frac{1}{i^2} \left[1+3(i-1)\cdot \frac{m-1}{N-1}+
(i-1)(i-2)\frac{(m-1)(m-2)}{(N-1)(N-2)}-((i-1)\frac{m-1}{N-1}+1)^2 \right].
\end{split}
\end{equation*}

Secondly, if $I_i=1$ the random variable $\operatorname{E}(P@i\cdot I_i \mid I_i)$ takes the value $\frac{1}{i}\cdot \big((i-1)\frac{m-1}{N-1}+1\big)$ with probability $\frac{m}{N}$ and $0$ otherwise. Therefore, 
\begin{equation*}
\begin{split} t_2&=\operatorname{Var}\big(\operatorname{E}(P@i\cdot I_i \mid I_i)\big)=\operatorname{E}\big((\operatorname{E}(P@i\cdot I_i \mid I_i))^2\big)-\big(\operatorname{E}(\operatorname{E}(P@i\cdot I_i \mid I_i))\big)^2=\\
&=\frac{1}{i^2}\left((i-1)\frac{m-1}{N-1}+1\right)^2\cdot \frac{m}{N}-\frac{1}{i^2}\left((i-1)\frac{m-1}{N-1}+1\right)^2\cdot \left(\frac{m}{N}\right)^2=\\
&=\frac{1}{i^2}\left((i-1)\frac{m-1}{N-1}+1\right)^2\cdot \frac{m}{N}\left(1-\frac{m}{N}\right).
\end{split}
\end{equation*}
So, finally, we have the following.
\begin{equation*}
\begin{split}
c_{ii}=&\operatorname{Var}\big(P@i\cdot I_i\big)=t_1+t_2=\\
= &\frac{m}{N} \cdot \frac{1}{i^2} \left(1+3(i-1)\cdot \frac{m-1}{N-1}+
(i-1)(i-2)\frac{(m-1)(m-2)}{(N-1)(N-2)}\right)-\\
&-\frac{m}{N} \cdot \frac{1}{i^2} \left((i-1)\frac{m-1}{N-1}+1\right)^2 +
\frac{1}{i^2}\left((i-1)\frac{m-1}{N-1}+1\right)^2\cdot \frac{m}{N}\left(1-\frac{m}{N}\right)=
\\
= &\frac{1}{i^2}\cdot \frac{m}{N} \left(1+3(i-1)\cdot \frac{m-1}{N-1}+
(i-1)(i-2)\frac{(m-1)(m-2)}{(N-1)(N-2)}-\right.\\
& -\left.\frac{m}{N}-2(i-1)\frac{m(m-1)}{N(N-1)} - 
 (i-1)^2\frac{m(m-1)^2}{N(N-1)^2} \right)=\\
= &\frac{1}{i^2}\cdot \frac{m}{N} \left(1- \frac{m}{N}+(i-1)\frac{m-1}{N-1}\left(3-2\frac{m}{N}+(i-2)\frac{m-2}{N-2}-(i-1)\frac{m}{N} \frac{m-1}{N-1}\right) \right)=\\
=& \frac{1}{i^2} \cdot \frac{m}{N} \left[1- \frac{m}{N}+\frac{m-1}{N-1} \left(2\frac{m-2}{N-2}-3+2\frac{m}{N}-\frac{m(m-1)}{N(N-1)}+ \right. \right.\\
&\left. \left. + i\left(3-2\frac{m}{N}-3\frac{m-2}{N-2}+2\frac{m(m-1)}{N(N-1)}\right) + i^2\left(\frac{m-2}{N-2}-\frac{m(m-1)}{N(N-1)}\right) \right) \right]=\\
= & \frac{m}{N} \left(1- \frac{m}{N}-\frac{m-1}{N-1} \left(3-2\frac{m-2}{N-2}-\frac{m}{N}\left(2-\frac{m-1}{N-1}\right)\right)\right)\frac{1}{i^2}+\\
& + \frac{m}{N}\frac{(m-1)}{(N-1)}\left(3\left(1-\frac{m-2}{N-2}\right)-2\frac{m}{N}\left(1-\frac{m-1}{N-1}\right)\right)\frac{1}{i}+\\ 
& + \frac{m}{N} \frac{(m-1)}{(N-1)}\frac{(m-2)}{(N-2)}-\left(\frac{m(m-1)}{N(N-1)}\right)^2 .
\end{split}
\end{equation*}
(ii)
For all such $i$ and $l$ that $1<i<l\leq k$
$$c_{il}=c_{li}=\operatorname{cov}(P@i\cdot I_i,P@l\cdot I_l)=\operatorname{E}(P@i\cdot I_i\cdot P@l\cdot I_l)-\operatorname{E}(P@i\cdot I_i)\cdot \operatorname{E}(P@l\cdot I_l).$$
Here, we can also use the conditional expectation for calculation
\[
\operatorname{E}(P@i I_i\cdot P@l I_l) = \operatorname{E}\big(\operatorname{E}(P@i I_i\cdot P@l I_l \mid I_i I_l)\big) 
= \operatorname{E}(P@i I_i\cdot P@l I_l \mid I_i I_l = 1)  P(I_i I_l = 1).
\]
If $I_iI_l=1$, then we know that the items in the $i$th and $l$th places are relevant, so we should take into account the possible arrangements of the other $m-2$ relevant items between the $N-2$ places. So,  
$P(I_iI_l=1)=\frac{C_{N-2}^{m-2}}{C_N^m}=\frac{m(m-1)}{N(N-1)}$. And then
\begin{equation*}
\begin{split}
\operatorname{E}(P@i\cdot & I_i\cdot P@l\cdot I_l \mid I_i I_l = 1)=\operatorname{E}\left(\frac{1}{i} \sum_{j=1}^{i} I_j 
\cdot \frac{1}{l} \sum_{r=1}^{l} I_r \mid I_i I_l = 1\right)=\\
= &\frac{1}{i\cdot l}\operatorname{E}\left(\left(1+\sum_{j=1}^{i-1}I_j\right)\left(1+\sum_{r=1}^{i-1}I_r+\sum_{r=i+1}^{l-1}I_r+1\right) \mid I_iI_l=1\right)=\\
= &\frac{1}{i\cdot l}\operatorname{E}\left(2+3\sum_{r=1}^{i-1}I_r+\sum_{r=i+1}^{l-1}I_r+\left(\sum_{j=1}^{i-1}I_j\right)^2+\sum_{j=1}^{i-1}I_j\sum_{r=i+1}^{l-1}I_r \mid I_iI_l=1\right)=\\
=&\frac{1}{i\cdot l}\operatorname{E}\left(2+4\sum_{r=1}^{i-1}I_r+\sum_{r=i+1}^{l-1}I_r+2\sum_{j=1}^{i-r}\sum_{r=j+1}^{i-1}I_jI_r+\sum_{j=1}^{i-1}\sum_{r=i+1}^{l-1}I_jI_r \mid I_iI_l=1\right)=\\
=&\frac{1}{i\cdot l}\left(2+4(i-1)\frac{m-2}{N-2}+(l-1-i)\frac{m-2}{N-2}+(i-1)(i-2)\frac{(m-2)(m-3)}{(N-2)(N-3)}+\right.\\
& \left. +(i-1)(l-1-i)\frac{(m-2)(m-3)}{(N-2)(N-3)}\right)=\\
= &\frac{1}{i\cdot l}\left(2+\frac{m-2}{N-2}(3i-5+l)+\frac{(m-2)(m-3)}{(N-2)(N-3)}(i-1)(l-3)\right).
\end{split}
\end{equation*}
Therefore, 
\begin{equation*}
\begin{split}
c_{i l} & = \frac{m(m-1)}{N(N-1)}\frac{1}{i\cdot l}\left[2-5\frac{m-2}{N-2}+i\cdot 3\frac{m-2}{N-2}+l\frac{m-2}{N-2}+i\cdot l\frac{(m-2)(m-3)}{(N-2)(N-3)}- \right.\\
&\left.-i\cdot 3\frac{(m-2)(m-3)}{(N-2)(N-3)}-l\frac{(m-2)(m-3)}{(N-2)(N-3)}+3\frac{(m-2)(m-3)}{(N-2)(N-3)}\right]-\\
&-\left(\frac{m}{N}\right)^2\frac{1}{i\cdot l}\left(1+(i-1)\frac{m-1}{N-1}\right)\left(1+(l-1)\frac{m-1}{N-1}\right)=\\
= & \frac{m}{N}\frac{1}{i\cdot l}\left[\frac{m-1}{N-1}\left(2-5\frac{m-2}{N-2}+3\frac{(m-2)(m-3)}{(N-2)(N-3)}+i\cdot 3\frac{m-2}{N-2}\left(1-\frac{m-3}{N-3}\right)+ \right.\right.\\
&\left. +l\frac{m-2}{N-2}\left(1-\frac{m-3}{N-3}\right)+i\cdot l\frac{(m-2)(m-3)}{(N-2)(N-3)}\right)-\frac{m}{N}\left(1+l\frac{m-1}{N-1}-\frac{m-1}{N-1}+ \right.\\
&+i\frac{m-1}{N-1}-\frac{m-1}{N-1} \left.\left.+il\left(\frac{m-1}{N-1}\right)^2-l\left(\frac{m-1}{N-1}\right)^2-i\left(\frac{m-1}{N-1}\right)^2+\left(\frac{m-1}{N-1}\right)^2\right)\right]=\\
= &\frac{m}{N}\frac{1}{i\cdot l}\left[\frac{m-1}{N-1}\left(2-5\frac{m-2}{N-2}+3\frac{(m-2)(m-3)}{(N-2)(N-3)}\right)-\frac{m}{N}\left(1-\frac{m-1}{N-1}\right)^2+ \right.\\
&+i\frac{m-1}{N-1}\left(3\frac{m-2}{N-2}\left(1-\frac{m-3}{N-3}\right)-\frac{m}{N}\left(1-\frac{m-1}{N-1}\right)\right)+\\
&+ l\frac{m-1}{N-1}\left(\frac{m-2}{N-2}\left(1-\frac{m-3}{N-3}\right)-\frac{m}{N}\left(1-\frac{m-1}{N-1}\right)\right)+\\
&+ \left. i\cdot l\frac{m-1}{N-1}\left(\frac{(m-2)(m-3)}{(N-2)(N-3)}-\frac{m}{N}\frac{m-1}{N-1}\right)\right].
\end{split}
\end{equation*}
Using the notations from the theorem's statement and  Lemma~\ref{lem1},  we'll get:
\begin{equation*}
\begin{split}
\operatorname{Var}&_{WOR}(AP@k)= \frac{1}{M^2} \Bigg( \sum_{i=1}^{k}c_{ii} + 2\sum_{i=1}^{k-1}\sum_{l=i+1}^{k}c_{il} \Bigg) =
\\
= &  \frac{1}{M^2} \Bigg[\sum_{i=1}^{k}\frac{m}{N}\Bigg(\frac{A}{i^2}+\frac{B}{i}+C \Bigg)+2\sum_{i=1}^{k-1}\sum_{l=i+1}^{k}\frac{m}{N}\Bigg(\frac{D}{i\cdot l}+\frac{E}{l}+\frac{F}{i}+G \Bigg) \Bigg] =
\\
= & \frac{1}{M^2} \frac{m}{N}\Bigg[ A \cdot H_k^{\left(2\right)} + B \cdot H_k + kC + D(H_k^2 - H_k^{\left(2\right)})  +\\
& +   2Ek(H_k-1) + 2F(k-H_k) + Gk(k-1) \Bigg] =
\\
= & \frac{1}{M^2} \frac{m}{N}\Bigg[ k[C+2(E-F) + (k-1)G] + H_k(B-2(E-kF))+\\
&+ H_k^2 \cdot D + H_k^{\left(2\right)}(A-D) \Bigg]. 
\end{split}
\end{equation*}
This completes the proof.
\end{proof}

\end{document}
